\documentclass[11pt]{article}

\usepackage[margin=0.82in]{geometry}
\usepackage{newtxtext,newtxmath}
\usepackage{microtype}
\usepackage{enumitem}
\usepackage{titlesec}
\usepackage[colorlinks=true,allcolors=blue]{hyperref}
\usepackage[numbers,sort&compress]{natbib}

\setlength{\parindent}{0pt}
\setlength{\parskip}{0.55em}
\titleformat{\section}{\large\bfseries}{}{0em}{}
\titlespacing*{\section}{0pt}{1.0em}{0.35em}

\begin{document}

\begin{center}
{\LARGE \textbf{Research Statement}}\\[0.35em]
{\large Yoshitaka Inoue}\\[0.25em]
\href{https://inoue0426.github.io/}{inoue0426.github.io}\\[0.25em]
{\small Last updated: October 2026}
\end{center}

\vspace{0.4em}

\section*{Learning how biological systems respond to intervention}

My research is organized around a simple question: \emph{what representation of a biological system is useful when the goal is to predict what will happen after we intervene?} In therapeutic discovery, this question appears at several scales. A model may need to predict whether a drug will work in a patient, how a cell state will change after treatment, which molecular mechanisms are likely to matter, or whether the evidence is too weak to support a confident conclusion. These problems are related, but they are often studied separately.

I want to connect them. My long-term goal is to develop machine-learning methods that learn how interventions transform biological systems and that transfer those learned transformations across experimental settings, modalities, and eventually patients. I am especially interested in the gap between the data that are abundant and experimentally controllable---cell lines, perturbation screens, molecular profiling---and the settings in which we ultimately want to make decisions, where measurements are sparse, heterogeneous, and confounded.

This perspective has emerged gradually from my work. My earlier research focused on drug-response prediction and interpretation. In \textbf{drGT}, published in \href{https://doi.org/10.1186/s12859-026-06417-z}{\emph{BMC Bioinformatics}}, I developed a heterogeneous drug--cell--gene graph model that couples response prediction with gene-level attribution \citep{inoue2026drgt}. In zero-shot transfer from GDSC1 to GDSC2, drGT reached $R^2=0.334$ and AUROC $=0.786$; 63.7\% of high-attention drug--gene links were supported by PubMed evidence or a structure-based drug--target model. The \href{https://github.com/sciluna/drGT}{code is public}.

That work raised a second problem: a model can propose a drug--gene relationship even when the surrounding evidence is incomplete or contradictory. I developed \textbf{DrugAgent} to integrate machine-learning, knowledge-graph, and literature evidence for drug--target interaction assessment \citep{inoue2026drugagent}. On a public benchmark of 900 kinase--drug pairs spanning 178 kinases and 42 inhibitors, its explanations were judged faithful to the supplied evidence in 98.8\% of cases, with 98\% label agreement across repeated runs. The \href{https://arxiv.org/abs/2408.13378}{preprint} and \href{https://github.com/sciluna/DrugAgent}{code} are public.

More recently, I shifted from static pretreatment prediction toward modeling the \emph{change induced by treatment}. In \textbf{PerturbRx}, we learn drug- and dose-conditioned latent transitions from single-cell perturbation data and transfer them to pretreatment patient profiles \citep{inoue2026perturbrx}. On TCGA-485, PerturbRx improved AUROC/AUPRC from $0.684/0.766$ for a Patient+Drug baseline to $0.709/0.794$; on the PDX benchmark, it improved from $0.673/0.587$ to $0.685/0.611$. The \href{https://arxiv.org/abs/2608.21349}{preprint is available on arXiv}.

These projects have changed how I think about therapeutic machine learning. The central object should not always be a static representation of a cell, tumor, or drug. In many settings, the more useful object is the \emph{relationship between a state, an intervention, and the state that follows}.

\section*{1. Intervention-conditioned representations}

Large perturbational datasets now make it possible to observe biological systems under many controlled interventions. Single-cell atlases such as Tahoe-100M contain responses to thousands of drug--context combinations at a scale that was difficult to imagine only a few years ago \citep{zhang2025tahoe}. Foundation models also provide increasingly expressive representations of cellular state \citep{hao2024scfoundation}. I use these resources to study reusable \emph{operators of change}: representations of how an intervention moves a biological system from one state to another.

Existing perturbation models have shown that latent-space arithmetic, factorized perturbation representations, and neural transport can capture aspects of cellular response and generalize across conditions \citep{lotfollahi2019scgen,lotfollahi2023cpa,bunne2023cellot}. A question I care about is what must be added for these representations to transfer beyond the experimental systems in which they were learned.

PerturbRx is one step in this direction. It uses the learned treatment-induced transition as a feature for patient-level drug response. The transition model is trained on controlled perturbation data, then frozen and applied to pretreatment patient samples. This separation is important to me. It asks whether a model can learn something about the \emph{effect of an intervention} that remains useful after the cell line, cohort, and prediction task have changed.

My next work in this area focuses on stricter forms of generalization. A model that performs well when the same compounds or cancer contexts are represented during training may still fail when both the intervention and biological background are unfamiliar. I therefore want to study generalization along explicit axes---unseen compounds, unseen biological contexts, and unseen compound--context combinations---rather than treating all held-out data as equivalent. I am also interested in identifying when a learned transition is genuinely transferable and when it is only a useful interpolation within the source domain.

A longer-term goal is to move from one-step treatment effects toward models of \emph{intervention-conditioned biological dynamics}. This includes dose, time, combinations, treatment sequence, and adaptation. The useful representation is one that supports decisions even when a post-treatment state is unavailable.

\section*{2. Which biological context should transfer?}

Transfer also depends on biological context. My work with spatial transcriptomics and cell--cell communication has made this concrete: a transcriptional profile can miss tissue organization and intercellular signals that shape treatment response.

My working hypothesis is that transferability depends on whether the source and target preserve the cell states and interactions required by a drug's mechanism of action (MoA). For one treatment, pathway activity within a tumor cell may carry most of the relevant signal. For another, response may depend on where that cell sits in tissue and which neighboring cell types are signaling to it.

I develop representations that keep these pieces of context explicit enough to test that hypothesis. Spatial structure, cell-state composition, and cell--cell communication can then be evaluated for their contribution to transfer, rather than being folded automatically into a single embedding. A failure to transfer is useful when it points to a missing biological context.

This gives patient stratification a more specific target: identify patients who preserve the biological context required for the same treatment mechanism to operate.

\section*{3. Mechanism-of-action, evidence, and knowing when not to predict}

Prediction in biology is more useful when the supporting evidence can be inspected. DrugAgent grew out of this problem. Biomedical evidence is fragmented across structured databases, predictive models, and experimental literature, and those sources frequently disagree.

My work treats mechanism-of-action (MoA) evidence as part of the prediction problem. A learned treatment transition can suggest a molecular direction of change. Target and pathway evidence can connect that change to a candidate MoA, while patient context determines whether the same explanation remains plausible in a new system. DrugAgent provides a framework for making disagreements among those evidence sources visible.

Abstention is the next part of this problem. Biomedical models are routinely evaluated outside the conditions represented during training. Calibrated uncertainty is useful, but I also want a model to identify what support is missing: an unfamiliar compound, weak target evidence, an unseen cell state, or a patient context that does not support the learned MoA.

In ongoing work, I am studying how missing evidence and distribution shift can be represented explicitly so that a model can withhold a prediction for a specific, inspectable reason.

\section*{A research program centered on transfer}

The common thread across these directions is transfer under intervention. I want to learn representations from experimental systems where interventions can be measured at scale, determine which biological structure makes those representations portable, and use mechanistic evidence to decide when transfer is justified.

This program combines representation learning, perturbation modeling, multimodal learning, causal questions, and scientific agents around one problem: building models that remain scientifically useful as we move from controlled experiments toward heterogeneous biological and clinical settings.

The practical benchmark I keep in mind is demanding. A model should learn from one set of interventions and contexts, make useful predictions in another, explain what evidence supports the transfer, and recognize when that evidence is too weak. That is the standard I use to connect experiments, mechanisms, and therapeutic decisions.

\bibliographystyle{unsrtnat}
\bibliography{references}

\end{document}
